
\section{\Увод}

Овај рад се бави {\sl логаритамском функцијом}, једном од најважнијих функција у математици.
Због своје важности, заједно са експоненцијалном,
тригонометријским и њима инверзним функцијама, спада у групу {\sl елементарних\/} функција. 
Описане су њене особине и дати примери њене употребе,
као и \hyperref[sec:zadaci]{задаци са решењама} (укупно \the\numexpr\totvalue{zadaci}).

Сама реч {\sl логаритам\/}\index{логаритам} потиче од грчких речи
$\lambda\acute o\gamma o\varsigma$~({\sl логос\/}) и 
$\alpha\rho\iota\theta\mu\acute o\varsigma$~({\sl аритмос\/}), 
са значењем \navod{одговарајући број}.



\subsection{Дефиниција логаритма}\index{дефиниција}

Функција
\begin{equation}
  \okvir{\scalebox{1.44}{$y=\log_b x$}}
\end{equation}
је решење по $y$ једначине
$$
x=b^{\,y},
$$
где је $b$ {\sl основа\/}\index{основа} ({\sl база\/})\index{база} логаритма, а $x$ {\sl\idx{аргумент}}.
(Изговара се \navod{$y$ је једнако логаритам од $x$ за основу $b$}
или краће \navod{$y$ је логаритам $b$ од $x$}.)


\subsection{Ток и график функције}

\newcounter{figno}
\def\newpic#1{\stepcounter{figno}%
  \hbox{{\bf Слика \thefigno:\/~}#1}}
\def\slika#1#2{\stepcounter{figno}\displaylines{
  \hbox{#1}\cr\hbox{{\bf Слика \thefigno:\/~}#2}}}

Функција је у скупу реалних бројева $\Rset$ дефинисана за $x>0$ и $b>0\land b\ne1$.
Функција је {\sl монотона\/}\index{монотоност}: за $b>1$ функција је {\sl растућа}, док за $b<1$ функција је {\sl опадајућа}.
Због тога важи {\sl бијекција\/}\index{бијекција}: $\logb u=\logb v \Leftrightarrow u=v$.
Функција има једну {\sl нулу}, увек за $x=1$. Када $x\to 0$, онда $y\to{-\infty}$\index{бесконачност $(\infty)$}
за $b>1$, односно, $y\to{+\infty}$ за $b<1$.
$$
\slika{\includegraphics[]{log.1}}{График\index{график} логаритамске функције $y=\logb x$.}
$$

\subsection{Антилогаритам}

\def\antilog{\mathop{\rm antilog}\nolimits}
Инверзна функција логаритму
је обично степеновање основе логаритма аргументом и зове се {\sl\idx{антилогаритам}}
\vskip-9pt
\begin{equation}
\okvir{\antilog_b x = \logb^{-1}x =  b^x}.
\end{equation}
Из саме дефиниције важи
\begin{equation}
\okvir{\log_b(\antilog_b x)=\antilog_b(\log_b x)=x}.
\end{equation}
